\section{Excluding forks and identifying Busemann pairs}\label{sec:forks}

The finite-tree count in this section follows the density-versus-boundary
principle of Burton and Keane~\cite{BurtonKeane1989}. Here the preceding
label budget bounds the edge--label pairs supplied by distinct branches.

Continue with the fixed chart $I$.  An \emph{upper fork at $x$} consists
of two bigeodesics, oriented toward upper labels in $I$, with the same
upper label, coincident backwards from $x$, and with different first
edges forwards from $x$.  Equivalently, one may require a common
backwards tail and a split at $x$: uniqueness of finite geodesics from a
vertex on that common tail identifies the two paths up to their split.
A lower fork is defined with the orientation reversed.  These are
existence events with two path witnesses and a label witness, and hence
are measurable in the completed space by \Cref{lem:measurable}.

\begin{proposition}[Absence of forks]\label{prop:no-forks}
Almost surely there is no upper or lower fork in any of the countably
many charts under consideration.
\end{proposition}

\begin{proof}
The counting argument groups forks by label and extracts disjoint
forward branches within each group. Positive fork density would force
a cubic expected number of usable edge--label pairs in large boxes;
uniform confinement and \Cref{prop:edge-count} give an arbitrarily small
cubic upper bound.

It suffices first to exclude upper forks in one chart.  Work on the
simultaneous event of \Cref{prop:opposite-coalescence}.  Every upper fork
label is exceptional by the first test in
\Cref{prop:exceptional-countable}, since two different rays of that label
start at its split vertex.  Consequently all lower rays with the
opposite label coalesce, simultaneously for every label arising in this
way.  In particular, a finite collection of bigeodesics with one such
upper label has a common backwards tail.  A vertex sufficiently far
along that common tail is a common past root for the entire finite
collection.

Let $F_n$ be the number of vertices in $[-n,n]^2\cap\Z^2$ that are upper
forks for this chart, and let
\[
 p_{\mathrm{fork}}=\PP(0\text{ is an upper fork for this chart}).
\]
The chart and the fork definition are translation invariant, so
\begin{equation}\label{eq:fork-count-mean}
 0\le F_n\le(2n+1)^2,
 \qquad \EE F_n=(2n+1)^2p_{\mathrm{fork}}.
\end{equation}
Fix $\eps>0$ and put $H_n=\lceil\eps n\rceil$.  Let $G_n$ be the
intersection of the event $U_n$ in \Cref{lem:usable} and the event that
$\mathcal R_e(H_n)$ holds for every edge with both endpoints in
$[-2n,2n]^2$.  Uniform all-ray confinement in \Cref{prop:cone} gives
the latter event eventually almost surely: all its endpoint sites lie in
a box of radius a fixed multiple of $n$, and their required offsets are
$o(n)$.  Therefore, for this fixed $\eps$,
\begin{equation}\label{eq:fork-good-event}
 \PP(G_n^c)\longrightarrow0.
\end{equation}
No rate uniform in $\eps$ is asserted or needed.

We claim the following deterministic counting inequality on $G_n$, for
$n\ge2$:
\begin{equation}\label{eq:fork-edge-credit}
 \frac n4 F_n
 \le\sum_{\substack{e:\text{both endpoints in }[-2n,2n]^2}}
          N_e(H_n).
\end{equation}
Choose a witnessing pair at each of the finitely many fork sites counted
by $F_n$, and group the choices according to their exact upper label.
These finite choices are only used to prove a pathwise inequality; no
measurable enumeration of labels is required.

Consider a group containing $q\ge1$ distinct fork sites.  There are at
most $2q$ witnessing bigeodesics in this group.  Their lower rays
coalesce by \Cref{prop:opposite-coalescence}.  Pairwise coalescence of a
finite collection gives a common tail by going beyond all their finitely
many merging locations.  Choose a vertex $o$ on this tail, earlier than
every selected split, and orient each witness forwards from $o$ toward
its upper end.

The union of these forward paths is a rooted tree.  To justify this
directly, the initial path from $o$ to any vertex $z$ on a witness is a
finite geodesic.  If two witnesses visit $z$, finite uniqueness makes
their entire paths from $o$ to $z$ equal.  They therefore cannot diverge
and later meet, and common edges cannot carry conflicting forward
orientations.  Each visited vertex other than $o$ has one parent, and
the resulting graph has no cycle.

Let $\mathcal T$ be the finite subtree obtained by taking the union of
the root paths to the two selected children immediately after every
selected split.  Each of the $q$ split vertices has at least two children
in $\mathcal T$.  For any finite rooted tree with $L$ leaves,
\begin{equation}\label{eq:finite-tree-leaves}
 L-1=\sum_{u:\,d^+_{\mathcal T}(u)>0}
               \bigl(d^+_{\mathcal T}(u)-1\bigr),
\end{equation}
where $d^+_{\mathcal T}(u)$ is its number of children.  Indeed, the sum of
the outdegrees is the number of edges, one less than the number of
vertices; subtracting the number of nonleaves gives
\eqref{eq:finite-tree-leaves}.  Thus $L\ge q+1$.

Every leaf of $\mathcal T$ is an endpoint of one of the root-to-child
paths defining $\mathcal T$.  Otherwise a defining path containing that
leaf would continue to another vertex of $\mathcal T$, contradicting its
being a leaf.  In particular, every leaf is a selected child of a fork
site, and hence lies in $[-n-1,n+1]^2$.  The root paths are allowed to
leave and reenter both boxes; only their terminal children determine this
location of the leaves.  Selected split vertices may also be ancestors
of one another, but each is still a distinct vertex with at least two
children and contributes at least one to
\eqref{eq:finite-tree-leaves}.

For each leaf choose a witnessing bigeodesic containing its defining
root-to-child path, and continue forwards along that witness.  These
continuations are pairwise vertex disjoint.  In fact, if two met at $w$,
uniqueness of the finite geodesic from $o$ to $w$ would put both leaves
on a single root path.  They would be comparable in the ancestor order,
contrary to being distinct leaves of $\mathcal T$.  The same reasoning
rules out an intersection between one continuation and the root path to
another leaf: its own leaf would then be an ancestor of the other one.
An intersection with an earlier point of its own root path contradicts
simplicity.  These conclusions hold throughout the plane, including
parts of $\mathcal T$ outside either box.

Each leaf has infinity-norm distance at least $n-1$ from the boundary of
$[-2n,2n]^2$.  Its continuation eventually reaches this boundary, since
an upper ray in the chart has height tending to $+\infty$.  Stop the
continuation at its first boundary visit.  The resulting path is simple,
all its edges have both endpoints in the outer box, and its length is at
least $n-1\ge n/2$.  Its starting point is in $[-3n,3n]^2$, so $U_n$
supplies at least half its edges as usable.  It therefore supplies at
least $n/4$ usable edges inside the outer box.  Every such edge lies on
the witnessing bigeodesic and satisfies $\mathcal R_e(H_n)$ on $G_n$.
It contributes the group's label to $N_e(H_n)$.

For this group the leaf continuations are disjoint, so no edge is
credited twice with that label.  Groups with different labels credit
different pairs $(e,\rho)$, even if they use the same geometric edge.
Thus the count is of distinct edge--label pairs.  Summing over the groups
credits at least
\[
 \frac n4\sum_{\text{groups}}L
 \ge\frac n4\sum_{\text{groups}}q
 =\frac n4 F_n
\]
pairs, proving \eqref{eq:fork-edge-credit}.

We now take expectations, retaining the scale of the bad-event loss.
By \eqref{eq:fork-count-mean},
\[
 \EE[F_n\mathbf 1_{G_n}]
 \ge (2n+1)^2
       \bigl(p_{\mathrm{fork}}-\PP(G_n^c)\bigr).
\]
The loss after multiplication by $n/4$ is at most
\begin{equation}\label{eq:fork-bad-event-loss}
 \frac n4(2n+1)^2\PP(G_n^c)=o(n^3)
 \qquad\text{for fixed }\eps.
\end{equation}
There are $2(4n)(4n+1)\le40n^2$ edges with both endpoints in the outer
box.  From \eqref{eq:fork-edge-credit} and \Cref{prop:edge-count},
\begin{equation}\label{eq:fork-expectation-bound}
 \frac n4(2n+1)^2
       \bigl(p_{\mathrm{fork}}-\PP(G_n^c)\bigr)
 \le 40C n^2\bigl(\lceil\eps n\rceil+1\bigr).
\end{equation}
Divide by $n^3$ and let $n\to\infty$ for this fixed $\eps$.  Equations
\eqref{eq:fork-good-event} and \eqref{eq:fork-expectation-bound} yield
$p_{\mathrm{fork}}\le40C\eps$.  Sending $\eps\downarrow0$ gives
$p_{\mathrm{fork}}=0$.  Countability of vertices excludes every upper
fork in this chart almost surely.  Reversing upper and lower gives the
same conclusion for lower forks, and the countable union over charts
proves the proposition.
\end{proof}

\begin{corollary}[Common Busemann pairs for a common label]
\label{cor:same-busemann}
On one event of probability one, whenever two oriented bigeodesics $p,q$
in any fixed chart have the same upper label, their two Busemann
functions agree:
\begin{equation}\label{eq:same-busemann-pair}
 B_{p^+}=B_{q^+},\qquad B_{p^-}=B_{q^-}.
\end{equation}
This assertion is simultaneous over the bigeodesics and their possibly
random labels.  It also holds almost surely in every environment whose
law is absolutely continuous with respect to the iid law, and
simultaneously for a countable collection of such environments.
\end{corollary}

\begin{proof}
Take the common full-measure event of
\Cref{prop:opposite-coalescence,prop:no-forks} and the structural
properties.  Let the two bigeodesics have upper label $\rho$.  If their
upper Busemann functions differ, $\rho$ passes the second exceptional
label test of \Cref{prop:exceptional-countable}.  The simultaneous
conclusion of \Cref{prop:opposite-coalescence} makes their lower rays
coalesce.  Choose a vertex $o$ sufficiently far along their common lower
tail.  The forward portions from $o$ are upper geodesic rays.  They
cannot have the same tail, since tail invariance in
\Cref{lem:busemann-basic} would make their upper Busemann functions
equal.  They must therefore have a first split.  Finite uniqueness
prevents any later meeting, so this split is an upper fork with a common
past, contrary to \Cref{prop:no-forks}.

If the lower Busemann functions differ, the lower label $-\rho$ is
exceptional by the same second test.  Apply
\Cref{prop:opposite-coalescence} with upper and lower reversed to
coalesce the upper rays.  Reversing the preceding argument then produces
a forbidden lower fork.  This proves both equalities in
\eqref{eq:same-busemann-pair}.

All these deductions are pathwise on the stated common event.  They do
not take an uncountable union of fixed-label null events, nor assert
coalescence of every one-sided family at random labels.  Absolute
continuity transfers the event to each individual changed environment;
countable intersection gives the final assertion.
\end{proof}
